For $s\in\mathscr F(U)$, write $\widetilde s:U\to\operatorname{Spe}(\mathscr F)$ for its germ map. The subsets $\widetilde s(U)$ are open in the specified topology: the inverse image under any other germ map is the set of points where the two germs agree, which is open. They cover the espace étalé and form a basis, since equality of germs implies equality of representatives on a smaller neighborhood. The projection restricts to a homeomorphism $\widetilde s(U)\to U$, with inverse $\widetilde s$.

A section of $\mathscr F^+$ is locally a germ map, so it is a continuous section of the projection. Conversely, let $t$ be a continuous section and choose a basic neighborhood $\widetilde s(V)$ of $t(P)$. On the open neighborhood $t^{-1}(\widetilde s(V))$ of $P$, the section condition forces $t(Q)=s_Q$. Thus $t$ is locally represented by sections of $\mathscr F$ and belongs to $\mathscr F^+$. These identifications commute with restriction. Consequently the canonical map identifies $\mathscr F(U)$ with all continuous sections for every $U$ exactly when $\mathscr F\cong\mathscr F^+$, that is, when $\mathscr F$ is a sheaf.
